\citet{Young:1954jz} include a collection of ``accounts given by old men to establish the area of the Navajo occupancy -- the Navajo domain'' under the title \textit{Traditional Navajo Country} (p. 10). The first account in this collection, which is presented here, is told by a Navajo medicine man, who recounts a story passed down by his grandfather, known as Man With A Moustache. Importantly, the text is fragmentary: \citet[13--14]{Young:1954jz} include an additional section in English that does not appear in the Navajo version. The narrative presented here includes only the translated portion of the original Navajo.

The narrative begins with the First People dwelling in the Underworld, where the mountains Sierra Blanca Peak (\textit{Sisnaajiní}), Mount Taylor (\textit{Tsoodził}), San Francisco Peak (\textit{Dook’o’oosłííd}) La Plata Mountain (\textit{Dibéntsaa}), Huerfano (\textit{Dził Ná’oodi\-łii}) and Gobernador Knob (\textit{Ch’óol’į́’í}) existed. Their story of emergence is marked by a disruption caused by Coyote, who captures the offspring of the Water Monster, instigating a series of events that force the people to seek escape as water rises in the Underworld. As they struggle, supernatural beings Sun and Moon aid them, and Locust eventually helps them find their way to the Earth's surface through La Plata Mountain. In their emergence, they bring with them elements essential to their survival, including seeds, plants, and knowledge of the Sacred Mountains, which define their new homeland. As the People establish themselves in this new world, First Man and First Woman position the four Sacred Mountains (\textit{Sisnaajiní, Tsoodził, Dook’o’ooslííd, Dibéntsaa}) in the cardinal directions. The storyteller emphasizes that these Sacred Mountains represent the heart of the Navajo world as both physical boundaries and spiritual anchors.
